\chapter*{Abstract (English)}
\addcontentsline{toc}{chapter}{Abstract (English)}

\noindent\textbf{Title:} Deep Learning-Based Early Warning System for Urban Flash Flooding in the Bangkok Metropolitan Area Using Distributed Sensor Networks and Satellite Rainfall Estimation\\[4pt]
\textbf{Author:} Kittipong Wattanasiri\\[4pt]
\textbf{Field of Study:} Water Resources Engineering\\[4pt]
\textbf{Thesis Advisor:} Assoc.\ Prof.\ Somsak Boonprasert, Ph.D.

\vspace{0.4cm}
\noindent
Flash flooding in the Bangkok Metropolitan Area (BMA) causes recurring disruption to transportation, commerce, and public safety, particularly during the southwest monsoon season. Existing warning systems rely primarily on rule-based thresholds applied to a sparse network of manual gauge stations, which yields lead times too short for effective evacuation and traffic rerouting. This thesis develops and evaluates a deep learning-based early warning framework that fuses ground-level water-level telemetry from a distributed Internet-of-Things (IoT) sensor network with satellite-derived rainfall estimates to forecast localized flash-flood risk with a lead time of up to three hours.

\noindent
A hybrid Convolutional Long Short-Term Memory (ConvLSTM) architecture was trained on 34 months of historical data collected from 42 water-level sensors deployed across six flood-prone districts, combined with half-hourly rainfall estimation products supplied under a data-sharing arrangement with Nexa Weather Analytics. The model was benchmarked against a physically based hydrodynamic model and two classical machine learning baselines (random forest and gradient-boosted trees). Results indicate that the proposed ConvLSTM framework achieves a root-mean-square error (RMSE) of 4.8~cm for 90-minute-ahead water-level forecasts, a 31\% improvement over the hydrodynamic baseline, while reducing false-alarm rate to 6.2\% at the operational alert threshold.

\noindent
The findings suggest that combining dense low-cost sensor telemetry with satellite rainfall products in a spatiotemporal deep learning model is a practical and cost-effective path toward extending flood-warning lead times in dense tropical urban environments. The thesis concludes with recommendations for phased deployment across the wider BMA drainage network and discusses limitations related to sensor maintenance and data continuity during severe storm events.

\vspace{0.4cm}
\noindent\textbf{Keywords:} flash flooding, early warning system, deep learning, ConvLSTM, IoT sensor networks, satellite rainfall estimation, Bangkok

% =====================================================================
% ABSTRACT (THAI, ROMANIZED)
% =====================================================================
